\section{Conclusion}
\label{sec:conc}

We presented AgentFail, a diagnostic framework that decomposes LLM-agent failures into five stages and distinguishes silent from loud failures. Across 5{,}984 runs (318 unique tasks, 5 models, 3 suites), we found that silent-failure rates are task-dependent (86\% on synthetic traps, 4\% on UCI, 19\% on DSBench), that agents never self-correct silent failures in the standard protocol (0\% recovery), but that feedback signals enable 15.7\% recovery overall---with silent failures remaining far harder to correct (10.4\%) than loud failures (36.7\%). A key diagnostic finding is that four of five taxonomy stages are instantiated in the data, while Code-Generation is genuinely absent (confirmed by both classifier and human annotation), indicating that current LLM agents do not make operation-level code errors on structured data-science tasks. Oracle-guided repair matches the no-op baseline (64\%), establishing an upper bound on reparability. We separate infrastructure failures (38.1\% of runs) from genuine model failures and provide an enriched manifest with 32 fields, validated by three human annotators (Fleiss' $\kappa = 0.860$; v2 classifier accuracy 76.6\% on gold). Future work should scale human annotation beyond 47 traces, generate tasks that trigger Code-Generation, and expand to code-interpreter and public agent frameworks.
